% RLJ main.tex Version 2026.1

\documentclass[10pt]{article} % For LaTeX2e

%%%%%%%%%%%%%%%%%%%%%%%%%%%%%%%%%%%%%%%%%%%%%%%%%%%%%%%%%%%%%%%%
% AUTHOR: Select ONE option:
%      [accepted]{rlj} --> for camera ready (after peer review, if accepted)
%      {rlj}           --> for submission
%      [preprint]{rlj} --> to de-anonymize and remove references to RLJ/RLC
%%%%%%%%%%%%%%%%%%%%%%%%%%%%%%%%%%%%%%%%%%%%%%%%%%%%%%%%%%%%%%%%
\usepackage{rlj}           % Should be uncommented for submission
%\usepackage[accepted]{rlj} % Should be uncommented for the camera-ready
%\usepackage[preprint]{rlj} % Should be uncommented for preprint versions

%%%%%%%%%%%%%%%%%%%%%%%%%%%%%%%%%%%%%%%%%%%%%%%%%%%%%%%%%%%%%%%%
% WARNING: The following packages are already included in the
%          rlj.sty style file:
%
%  1. fancyhdr  - For controlling header/footers
%  2. natbib    - For formatting the bibliography
%  3. enumitem  - To customize the appearance of lists
%  4. fontenc (with option [T1]) - Allows for proper hyphenation and accents
%  5. times     - Times new roman font
%  6. ragged2e  - Used to justify text
%  7. tcolorbox - Used to create boxes on cover page
%  8. hyperref  - Configures hyperlinks throughout (e.g., links to references)
%  9. xcolor    - Used to define custom colors for links and boxes
%  10. amsmath  - Not used, but conflicts with lineno, so we include (and patch) it for authors
%  11. etoolbox - Included in the amsmath + lineno patch
%  12. lineno   - For adding line numbers when in submission
%
% You do not need to include them again in your main.tex.
% Including them again may lead to conflicts or compilation errors.
% Additionally, avoid loading packages that might conflict with these.
%%%%%%%%%%%%%%%%%%%%%%%%%%%%%%%%%%%%%%%%%%%%%%%%%%%%%%%%%%%%%%%%

%%%%%%%%%%%%%%%%%%%%%%%%%%%%%%%%%%%%%%%%%%%%%%%%%%%%%%%%%%%%%%%%
% Recommended (but not required) packages
%%%%%%%%%%%%%%%%%%%%%%%%%%%%%%%%%%%%%%%%%%%%%%%%%%%%%%%%%%%%%%%%
\usepackage{amsthm}
\usepackage{amssymb}            % Defines common symbols like \mathbb R
\usepackage{mathtools}          % Extends amsmath, providing common math tools
\usepackage{mathrsfs}           % Enables \mathscr, which can work in cases that \mathcal does not
%\mathtoolsset{showonlyrefs}     % Only number equations that are referenced (optional)
\usepackage{graphicx}           % For including images
\usepackage{subcaption}         % Allows for the use of subfigures and subcaptions
\usepackage[space]{grffile}     % For spaces in image names
\usepackage{url}                % For displaying URLs
\usepackage{lipsum}             % For placeholder text
\usepackage{tikz}
\usepackage{svg}
\usepackage{algorithm2e}
\usepackage{booktabs} % for professional tables
\usepackage[english]{babel}
\theoremstyle{definition}
\newtheorem{definition}{Definition}
%%%%%%%%%%%%%%%%%%%%%%%%%%%%%%%%%%%%%%%%%%%%%%%%%%%%%%%%%%%%%%%%
% AUTHOR: Fill in the following meta-data
%%%%%%%%%%%%%%%%%%%%%%%%%%%%%%%%%%%%%%%%%%%%%%%%%%%%%%%%%%%%%%%%

% Enter the title of your paper:
\title{Limits of reinforcement learning for decision trees in\\ Markov decision processes}

% The "running title" will be displayed in the header on every-other page.
% It is typically either the same as the title or a shorter version of the title.
% Enter your running title here:
\setrunningtitle{Limits of RL for decision tree policies}

% WARNING: Authors must not appear in the submitted version. They should be hidden
% as long as the rlj package is used without the [accepted] or [preprint] options.
% Non-anonymous submissions will be rejected without review.

% Enter the author names below. 
% NOTE: Denote affiliations using superscripts as in the provided example.
% NOTE: Use \textscript{1,2,3} instead of $^{1,2,3}$.
%       - Failure to do so will cause affiliation superscripts to appear on the cover page for camera-ready and preprint versions.
\author{Hector Kohler\textsuperscript{4,$\dagger$},
Riad Akrour\textsuperscript{2,1,3},
Philippe Preux\textsuperscript{1,2,3},}

% NOTE: For camera-ready and preprint versions, the cover page includes author names but not affiliations.
% It automatically removes the superscripts for affiliations.
% If the automatic process breaks (e.g., if an author name should include a superscript), you can manually define the author string to appear on the cover page by uncommenting the following line.
%\coverPageAuthor{Marlos C. Machado, Philip S. Thomas, Lorem Ipsum}

% Author emails, which can be clustered if they have shared endings as in this example
\emails{hector.kohler@polytechnique.edu, \ 
\{riad.akrour,philippe.preux\}@inria.fr}

% Author affiliations, in the order the occur
% The inclusion of state/province, etc. is optional.
% The inclusion of multiple affiliations is optional.
%   - List multiple affiliations with comma-separated numbers as in the example.
\affiliations{
$^{1}$\textbf{Universit\'e de Lille, France}\\
$^{2}$\textbf{Inria, France}\\
$^{3}$\textbf{UMR 9198-CRIStAL, CNRS, Centrale Lille, France}\\
$^{4}$\textbf{LIX, Ecole polytechnique, CNRS, Institut Polytechnique de Paris, France}\\
$\dagger$ Corresponding author}

\contribution{
    % Contribution
    We bridge the gap between RL for decision tree policies and POMDPs.
    }
    {
    \citet{topin2021iterative} proposed the first RL algorithms that train decision tree policies to maximize the rewards in MDPs without having to first compute an expert policy, e.g. a neural network.
    Contemporarily, \citet{pinto,baisero-ppo,baisero-dqn} developed asymmetric RL algorithms that are specifically tailored to train policies in POMDPs. We unify those two approaches.
    }

\contribution{
    % Contribution
    We show that RL for decision tree policies is hard because of partial observability.
    }
    {
    We perform comprehensive empirical benchmarking of asymmetric and standard RL on simple decision tree policy learning tasks for which we know the optimal solutions and show that RL fails. We relate our findings to hardness results from the POMDP literature~\cite{littman1}.
    Furthermore, we show that RL can consistently learn optimal decision tree policies for task that are fully observable which further corroborates our findings.
    }

% Include a list of keywords for the topic of the paper:
\keywords{Reinforcement learning, decision trees, MDPs, POMDPs} % Your keywords

% Define the summary that appears on the cover page.
\summary{For applications like medicine, machine learning models ought to be interpretable.
In that case, decision tree models are preferred over neural networks because humans can read their predictions from the root to the leaves.
Training such decision trees for sequential decision making problems is a relatively new research direction and most of the existing literature focuses on imitating neural networks.
In contrast, we study reinforcement learning (RL) algorithms that train decision trees \textit{directly} optimizing some trade-off of cumulative rewards and interpretability in a Markov decision process (MDP). 
We show that such algorithms can be seen as training policies for partially observable Markov decision processes (POMDPs).
This helps us understand why in practice it is often hard to train decision tree policies from scratch in MDPs.}

%%%%%%%%%%%%%%%%%%%%%%%%%%%%%%%%%%%%%%%%%%%%%%%%%%%%%%%%%%%%%%%%
%% Begin document, create title and abstract
%%%%%%%%%%%%%%%%%%%%%%%%%%%%%%%%%%%%%%%%%%%%%%%%%%%%%%%%%%%%%%%%
\begin{document}

\makeCover  % Create the cover page
\maketitle  % Make the title section

\begin{abstract}
For applications like medicine, machine learning models ought to be interpretable.
In that case, decision tree models are preferred over neural networks because humans can read their predictions from the root to the leaves.
Training such decision trees for sequential decision making problems is a relatively new research direction and most of the existing literature focuses on imitating neural networks.
In contrast, we study reinforcement learning (RL) algorithms that train decision trees \textit{directly} optimizing some trade-off of cumulative rewards and interpretability in a Markov decision process (MDP). 
We show that such algorithms can be seen as training policies for partially observable Markov decision processes (POMDPs).
This helps us understand why in practice it is often hard to train decision tree policies from scratch in MDPs.
\end{abstract}

%%%%%%%%%%%%%%%%%%%%%%%%%%%%%%%%%%%%%%%%%%%%%%%%%%%%%%%%%%%%%%%%
%% Section: Submission of papers to RLJ/RLC
%%%%%%%%%%%%%%%%%%%%%%%%%%%%%%%%%%%%%%%%%%%%%%%%%%%%%%%%%%%%%%%%
\section{Introduction}
Interpretability in machine learning is commonly divided into local and global approaches~\citep{glanois-survey}.
Local methods—also referred to as explainability or post-hoc methods~\citep{lipton}—provide explanations for individual predictions using tools such as local linear approximations~\citep{lime}, saliency maps~\citep{Puri2020Explain}, feature attributions~\citep{shap}, or attention mechanisms~\citep{attention}.
Although widely used, these methods approximate the behavior of an underlying black-box model and may therefore be unfaithful to its true computations~\citep{Atrey2020Exploratory}.

Global interpretability approaches instead restrict the model class so that the trained model is transparent by construction.
Decision trees~\citep{breiman1984classification} are a canonical example, as their predictions can be inspected, reasoned about, and formally verified.
This makes them particularly attractive for safety-critical applications and has motivated extensive research in supervised learning~\citep{lookahead,binoct,murtree,blossom,pystreed}.

Extending global interpretability to sequential decision making, however, remains challenging.
Existing approaches largely rely on \emph{indirect} methods~\citep{milani-survey}: a high-performing but opaque policy (typically a neural network) is first trained using reinforcement learning, and an interpretable model is then trained to imitate its behavior.
A prominent example is VIPER~\citep{viper}, which distills neural network policies into decision trees using imitation learning~\citep{dagger}.
Such methods have demonstrated strong empirical performance and enable formal verification~\citep{maraboupy}, but they optimize a surrogate objective—policy imitation—rather than the original reinforcement learning objective.
As a result, the best decision tree policy for the task may differ substantially from the tree that best approximates a neural expert. The curious reader will find an example of this phenomenon in appendix~\ref{sec:imit}.

This limitation motivates the study of \emph{direct} approaches that train interpretable policies by optimizing the reinforcement learning objective itself.
While direct decision tree learning is well understood in supervised settings, it is far less developed for sequential decision making.
Understanding why direct optimization is hard—and when it can succeed—is the central focus of this work.

In this article, we show that reinforcement learning of decision tree policies for MDPs, i.e. learning a decision tree that directly optimizes the cumulative reward of the process without relying on a black-box expert, is often very hard.
To do so, we construct very simple MDPs for which we know optimal decision tree policies and show that RL consistently fails to retrieve those policies.
We identify partial observability as a key reason for those failures. 

In section~\ref{sec:related-work-pomdp}, we present the related work on reinforcement learning algorithms that train decision tree policies for MDPs.
In section~\ref{sec:prelims}, we present key concepts for decsion trees, MDPs, and the formalism from~\citet{topin2021iterative} for reinforcement learning of decision tree policies.
In section~\ref{sec:poibmdp}, we show that this direct approach is equivalent to learning a \emph{deterministic memoryless} policy for POMDPs~\citep{POMDP}.
In section~\ref{sec:methodology}, we present our methodology to benchmark RL algorithms that train decision tree policies.
In section~\ref{sec:experiments}, we show that when RL fails to retrieve optimal decision tree policies for MDPs it is most likely because partial observability is involved.

\section{Related work}\label{sec:related-work-pomdp}
There exist reinforcement learning algorithms that directly train decision tree policies optimizing the cumulative rewards in a given MDP.
These approaches can be divided into methods based on \emph{parametric} and \emph{non-parametric} trees.

Parametric decision trees fix the tree structure \emph{a priori}—the depth and node arrangements—and only learn the node labels.
This formulation enables differentiability and allows direct optimization of the MDP cumulative rewards using policy gradient methods~\citep{pg_sutton}.
Several works employ PPO to train such differentiable trees~\citep{ppo,silva,vos2024optimizinginterpretabledecisiontree,sympol}.
While these methods can achieve strong performance, they require the tree structure to be specified in advance, making it difficult to adaptively trade off interpretability and performance.
An overly complex structure may require post-hoc pruning, whereas an insufficiently expressive structure may fail to represent good policies.
Moreover,~\cite{sympol} reports that additional stabilization techniques, such as adaptive batch sizes, are often necessary for direct reinforcement learning to match indirect imitation methods performances.

Non-parametric decision trees, by contrast, are the standard model in supervised learning, where algorithms efficiently construct trees that balance predictive performance and interpretability~\citep{breiman1984classification,oct}.
This trade-off is optimized by using a regularized training objective. However, training non-parametric trees to trade-off MDP cumulative rewards and interpretability is largely unexplored~\citep{milani-survey}.
To the best of our knowledge, the only work that studies this setting is~\cite{topin2021iterative}.
\cite{topin2021iterative} introduce \emph{iterative bounding MDPs} (IBMDPs), which augment the downstream MDP with additional state features, actions, rewards, and transitions.
They show that certain policies in the IBMDP correspond to decision tree policies for the downstream MDP. Hence, RL algorithms such as Deep Q-Networks~\citep{dqn} can be used to learn such policies in IBMDPs.

Finally, a few specialized methods exist for restricted problem classes. For maze-like MDPs,~\cite{theory1} proves the existence of optimal decision tree policies and provides a constructive algorithm.
In settings where the MDP model is fully known,~\cite{dt-opt-mdp} use planning to compute shallow parametric decision tree policies. Next, we recall useful technical material.
\input{technicals}
\section{Bridging the gap with the partially observable MDPs literature}\label{sec:poibmdp}
To better understand RL of decision tree policies for MDPs, we formulate the problem of training a deterministic policy depending on current observations in an IBMDP as training a deterministic policy depending only on current observations--also known as a deterministic \emph{memoryless} policy--in a partially observable Markov decision process (POMDP,~\cite{POMDP,chap2}).
By doing so, we can leverage results from the POMDP literature that is richer than interpretable RL literature.

\subsection{An adequate formalism}
\begin{definition}[Partially observable Markov decision process]\label{def:pomdp}
A partially observable Markov decision process is a tuple $\langle X, A, O, R, T, T_0, \Omega\rangle$. $X$ is the hidden state space.
$A$ is a finite set of actions.
$O$ is a set of observations.
$T: X \times A \rightarrow \Delta(X)$ is the transition function, where $T(\boldsymbol{x}_t, a, \boldsymbol{x}_{t+1}) = P(\boldsymbol{x}_t|\boldsymbol{x}_{t+1}, a)$ is the probability of transitioning to state $\boldsymbol{x}_{t}$ when taking action $a$ in state $\boldsymbol{x}$.
$T_0$: is the initial distribution over states. 
$\Omega: X \rightarrow \Delta(O)$ is the observation function, where $\Omega(\boldsymbol{o}, a, \boldsymbol{x}) = P(\boldsymbol{o}|\boldsymbol{x}, a)$ is the probability of observing $\boldsymbol{o}$ in state $\boldsymbol{x}$.
$R: X \times A \rightarrow \mathbb{R}$ is the reward function, where $R(\boldsymbol{x}, a)$ is the immediate reward for taking action $a$ in state $\boldsymbol{x}$.
Note that $\langle X, A, R, T, T_0 \rangle$ defines an MDP (definition~\ref{def:mdp}).
\end{definition}

Next, we can define partially observable iterative bounding Markov decision processes (POIBMDPs). 
They are IBMDPs (definition~\ref{def:ibmdp}) for which we explicitly define an observation space and an observation function.

\begin{definition}[Partially observable iterative bounding Markov decision process]\label{def:poibmdp} a partially observable iterative bounding Markov decision process $\mathcal{M}_{POIB}$ is a tuple:
    \begin{align*}
        \langle \overbrace{S\times O}^{\text{States}}, \overbrace{A\cup A_{info}}^{\text{Action space}},\overbrace{O}^{\text{Observations}} ,\overbrace{(R, \zeta)}^{\text{Rewards}}, \overbrace{(T_{info}, T, T_0)}^{\text{Transitions}}, \Omega \rangle
    \end{align*}
Where $\langle S\times O, A\cup A_{info}, (R, \zeta),( T, T_0, T_{info})\rangle$ is an IBMDP (definition~\ref{def:ibmdp}).
The transition function $\Omega$ maps concatenation of state features and observations--IBMDP states--to observations, $\Omega:S\times O \rightarrow O$, with $P(\boldsymbol{o}|(\boldsymbol{s}, \boldsymbol{o}))=1$ 
\end{definition}
POIBMDPs are particular instances of POMDPs (definition~\ref{def:pomdp}) where the observation function simply applies a mask over some features of the hidden state.
This setting has other names in the literature.
For example, POIBMDPs are mixed observability MDPs~\citep{momdp} with downstream MDP state features as the \textit{hidden variables} and feature bounds as \textit{visible} variables.
POIBMDPs can also be seen as non-stationary MDPs (N-MDPs)~\citep{learning-pomdp} in which there is one different transition function per downstream MDP state: these are called hidden-mode MDPs~\citep{hmmdp}.
Following~\cite{learning-pomdp}, we can write the value of a deterministic memoryless policy $\pi:O\rightarrow A\cup A_{info}$ in observation $\boldsymbol{o}$.

\begin{definition}[Value of an observation]\label{def:vpo} In a POIBMDP (definition~\ref{def:poibmdp}), the expected cumulative discounted reward of a deterministic memoryless policy $\pi:O\rightarrow A\cup A_{info}$ starting from observation $o$ is $V^{\pi}(\boldsymbol{o})$:
$V^{\pi}(\boldsymbol{o}) = \underset{(s,\boldsymbol{o}')\in S\times O}{\sum}P^{\pi}((\boldsymbol{s}, \boldsymbol{o}')|\boldsymbol{o})V^{\pi}((\boldsymbol{s}, \boldsymbol{o}'))$
With $P^{\pi}((\boldsymbol{s}, \boldsymbol{o}')|\boldsymbol{o})$ the asymptotic occupancy distribution (see section 4 from~\cite{learning-pomdp} for the full definition) of the hidden POIBMDP state $(\boldsymbol{s},\boldsymbol{o}')$ given the partial observation $o$ and $V^{\pi}((s, o'))$ the classical state-value function (definition~\ref{def:vs}).
We abuse notation and denote both values of observations and values of states by $V$ since the function input is not ambiguous.
\end{definition}

\subsection{Reinforcement learning in POMDPs}\label{sec:asym}

In general, the policy that maximizes the RL objective (definition~\ref{def:mdp-obj}) in a POMDP maps ``belief states'' or observation histories to actions~\citep{chap2}.
Hence, those policies do not correspond to decicion trees since we require that policies depend only on the current observation (algorithm~\ref{alg:extract-tree}).
If we did not have this constraint, we could apply any standard RL algorithm to solve POIBMDPs by seeking policies depending on belief states or observations histories because those are sufficient to optimally control any POMDP~\cite{chap2,lambrechts2025informed}.

In particular, the problem of finding the optimal deterministic memoryless policies for POMDPs is NP-HARD, even with full knowledge of transitions and rewards(section 3.2 from~\cite{littman1}).
It means that it is impractical to enumerate all possible policies and take the best one. 
For even moderate-sized POMDPs, a brute-force approach would take a very long time since there are $|A|^{|O|}$ deterministic memoryless policies.
Hence it is interesting to study RL for training deterministic memoryless policies since it would not enumerate the whole solution space.

In~\cite{learning-pomdp}, the authors show that the optimal memoryless policy can be stochastic. Hence, policy gradient algorithms~\citep{pg_sutton}--that return stochastic policies--are to avoid since we seek the best \textit{deterministic} policy. 
Furthermore, the optimal deterministic memoryless policy might not maximize all the values of all observations simultaneously~\citep{learning-pomdp} which makes it difficult to use TD-learning algorithms like Q-learning~\citep{watkins1992q}.
Indeed, doing a TD-learning update of one observation's value can change the value of \textit{all} other observations in an uncontrollable manner because of the dependence in $P^{\pi}((s, \boldsymbol{o}')|\boldsymbol{o})$ (definition~\ref{def:vpo}).

\RestyleAlgo{ruled}
\SetKwComment{Comment}{}{}
\begin{algorithm}
    \KwData{A POMDP, learning rates $\alpha_u,\ \alpha_q$, exploration prob. $\epsilon$}
    \KwResult{$\pi:O\rightarrow A$}
    \textcolor{teal}{Initialize $U(\boldsymbol{x},a) = 0$ for all $x \in X, a \in A$ \\}
    Initialize $Q(\boldsymbol{o},a) = 0$ for all $o \in O, a \in A$ \\
    \For{each episode}{
        Initialize state $x_0 \sim T_0$ \\
        Initialize observation $\boldsymbol{o}_0 \sim \Omega(\boldsymbol{x}_0)$ \\

        \For{each step $t$}{
            Choose action $a_t$ using $\epsilon$-greedy: $a_t = \operatorname{argmax}_a Q(\boldsymbol{o}_t,a)$ with prob. $1-\epsilon$ \\
            Take action $a_t$, observe $r_t = R(\boldsymbol{x}_t,a_t)$, $x_{t+1} \sim T(x_t,a_t)$, and $\boldsymbol{o}_{t+1} \sim \Omega(\boldsymbol{x}_{t+1})$ \\
            \textcolor{teal}{$y \leftarrow r + \gamma U(\boldsymbol{x}_{t+1}, \operatorname{argmax}_{a'} Q(\boldsymbol{o}_{t+1}, a'))$} \\
            $U(\boldsymbol{x}_t,a_t) \leftarrow (1 - \alpha_u) U(\boldsymbol{x}_t, a_t) + \alpha_u y $ \\
            \textcolor{teal}{$Q(\boldsymbol{o}_t,a_t) \leftarrow (1 - \alpha_q) Q(\boldsymbol{o}_t, a_t) + \alpha_q y $} \\
            $x_t \leftarrow \boldsymbol{x}_{t+1}$ \\
            $\boldsymbol{o}_t \leftarrow \boldsymbol{o}_{t+1}$ \\
        }
    }
    $\pi(o) = \operatorname{argmax}_a Q(\boldsymbol{o},a)$
    \caption{Asymmetric Q-Learning~\citep{baisero-dqn}. We highlight in green the differences with the standard Q-learning~\citep{watkins1992q}.}\label{alg:asymqlearning}
\end{algorithm}

Despite those hardness results, applying RL to POMDPs, by naively replacing the processes (hidden) states $\boldsymbol{x}$ by the observation $\boldsymbol{o}$ in Q-learning or Sarsa~\citep{sutton}, has already demonstrated successful in practice~\citep{sarsa-pomdp}. More recently, the framework of asymmetric RL~\citep{pinto,baisero-dqn,baisero-ppo} has shown promising results to learn policies for POMDPs. Asymmetric RL algorithms train a model--a policy or a value function--depending on hidden state (only available at train time) and a history dependent (or observation dependent) model to be deployed at test time.
The history or observation dependent model serves as target or critic to train the hidden state dependent model.
The history dependent (or observation dependent) model can thus be deployed in the POMDP after training since it does not require access to the hidden state to output actions.
In appendix~\ref{sec:hp-pomdp}, we present asymmetric variants of standard tabular RL algorithm.
For example, asymmetric Q-learning (algorithm~\ref{alg:asymqlearning})is a variant of Q-learning that returns a deterministic memoryless policy.
Given a POMDP, asymmetric Q-learning trains a partially observable Q-function $Q:O\times A\rightarrow\mathbb{R}$ and a Q-function $U:X\times A\rightarrow\mathbb{R}$.
The hidden-state-dependent Q-function $U$ serves as a target in the temporal difference learning update.
It is the tabular version of the modified DQN algoritm used in~\cite{topin2021iterative}.
Indeed, when learning deterministic memoryless policies for IBMDPs,~\cite{topin2021iterative} were using RL algorithms corresponding to asymmetric DQN or asymmetric PPO before those were published~\citep{baisero-dqn,baisero-ppo}.
We provide a reproducibility study of~\cite{topin2021iterative} in appendix~\ref{sec:reprod}.
Until recently, the benefits of asymmetric RL over standard RL was only shown empirically and only for history-dependent models.
The work of~\cite{justif-asym} proves that some asymmetric RL algorithms should in theory learn better history-dependent \textbf{or} memoryless policies for POMDPs which is exactly what we wish for.
However, those algorithms are inpractical because they require estimations of the asymptotic occupancy distribution $P^{\pi}((s, \boldsymbol{o}')|\boldsymbol{o})$ (definition~\ref{def:vpo}) for candidate policies which in turn would require to perform costly Monte-Carlo estimations.
We leave it to future work to use those algorithms that combine asymmetric RL and estimation of future visitation frequencies since those results are contemporary to the writing of this manuscript.

Brdiging the gap between RL of decision tree policies in MDPs through the formalism of IBMDPs (sections~\ref{sec:ibmdp} and~\ref{sec:extract-tree}) and RL of deterministic memoryless policies for POMDPs gave us insights on the hardness fo the task.
More importantly, it gave us insights on what RL algorithms could be used to train decision tree policies.
In the next section, we benchmark those asymmetric RL algorithms to train decision tree policies (definition~\ref{sec:asym}).
\input{rl_pomdp}

\section{Discussion}
In this paper, we were interested in algorithms that train decision tree policies to directly optimize some trade-off of interpretability and performance in MDPs without relying on an oracle or expert policy.
Starting from the IBMDP formulation from~\cite{topin2021iterative}, we have shown that such algorithms essentially train deterministic memoryless policies in POMDPs that we called POIBMDPs.
By bridging the gap with the POMDP literature, we gave strong insights on the hardness of this task as well as on existing RL algorithms specifically tailed for learning in POMDPs.
We supported those insights by benchmarking those RL algorithms on carefully crafted problems for which we knew the exact optimal decision tree policies.
Across our experiments, we found that only when partial observablity is absent from the task can good decision tree policies be trained.
Attempting to overcome the partial observability challenges seems like a bad research direction.
Indeed, while algorithms tailored specifically for the problem of training deterministic memoryless policies for POIBMDPs might exist, imitation learning works well in practice and has been the state-of-the-art for interpretable sequential decision making for a while.
Furthermore there are other limitations to IBMDPs that we did not cover such as how to choose good candidates information-gathering actions when building the decision tree policy or simply how to choose $\zeta$ for a target trade-off between tree interpretablility and downstream performance.
% Finally, while we focused on non-parametric tree learning assuming RL algorithms should naturally trade off interpretability and performance through the reward signal $\zeta$ for adding nodes to the decision tree policy, another future research avenue could be to develop better algortihms training parametric trees.
% Indeed parametric tree policies, on the other hand, can be computed with reinforcement learning directly in the downstream MDP.
% However such RL algorithms for parametric decision tree policies~\cite{silva,vos2024optimizinginterpretabledecisiontree,sympol} require to re-train a policy entirely for each desired level of interpretability, i.e. each unique tree structure.
% Future research in this direction should focus on algorithms for parametric tree policies that can re-use samples from one tree learning to train a different tree structure more efficiently.
% This would reduce the required quantity of a priori knowledge on the decision tree policy structure.
% \subsubsection*{Broader Impact Statement}
% \label{sec:broaderImpact}
% Our work put great emphasis on covering existing work, open sourcing code, and being transparent about limitations.
% We hope that the impact of our work is going to be positive for society through advancing reseach in interpretable machine learning.
% \subsubsection*{Acknowledgments}
% \label{sec:ack}
% Use unnumbered third level headings for the acknowledgments. All acknowledgments, including those to funding agencies, go at the end of the paper. Only add this information once your submission is accepted and deanonymized. The acknowledgments do not count towards the 8--12 page limit.

%%%%%%%%%%%%%%%%%%%%%%%%%%%%%%%%%%%%%%%%%%%%%%%%%%%%%%%%%%%%%%%%
%% NOTE: THIS MARKS THE END OF THE "MAIN TEXT"
%%%%%%%%%%%%%%%%%%%%%%%%%%%%%%%%%%%%%%%%%%%%%%%%%%%%%%%%%%%%%%%%

%%%%%%%%%%%%%%%%%%%%%%%%%%%%%%%%%%%%%%%%%%%%%%%%%%%%%%%%%%%%%%%%
%% Bibliography
%%%%%%%%%%%%%%%%%%%%%%%%%%%%%%%%%%%%%%%%%%%%%%%%%%%%%%%%%%%%%%%%
\bibliography{main}
\bibliographystyle{rlj}

%%%%%%%%%%%%%%%%%%%%%%%%%%%%%%%%%%%%%%%%%%%%%%%%%%%%%%%%%%%%%%%%
% AUTHOR: If your paper has no supplementary materials, you may 
%         comment out the line below, which creates the title for
%         the supplementary materials.
%%%%%%%%%%%%%%%%%%%%%%%%%%%%%%%%%%%%%%%%%%%%%%%%%%%%%%%%%%%%%%%%
\beginSupplementaryMaterials
\input{appendix}
\end{document}
